\chapter{행렬식}

\chapterintro{행렬식은 정사각행렬에 부여되는 스칼라 값으로, 가역성 판정과 선형변환의 부피 변화, 그리고 이후의 고유값 이론을 하나로 묶어 주는 핵심 도구입니다. 이 장에서는 계산 공식만 외우는 방식이 아니라, 왜 그 공식이 성립하는지를 교대 다중선형성 관점에서 차근차근 확인하겠습니다.}

\chaptergoals{Leibniz 공식으로 행렬식을 정의하고, 존재와 유일성을 스스로 설명할 수 있다.}{행연산, 곱셈성, 여인수 전개, 수반행렬 공식을 엄밀한 증명과 함께 사용할 수 있다.}{Cramer 공식과 Vandermonde 행렬식을 통해 행렬식의 계산적/이론적 응용을 연결할 수 있다.}

\sectiontemplate{Leibniz 공식과 행렬식의 존재/유일성}{행렬식은 임의로 정해진 함수가 아니라, 교대성과 다중선형성 그리고 정규화 조건으로 사실상 유일하게 결정됩니다. 먼저 순열 합으로 주어지는 Leibniz 공식(Leibniz formula)을 기준점으로 삼겠습니다.}{\item 순열의 부호와 Leibniz 공식을 정확히 쓸 수 있습니다.\item Leibniz 공식이 교대 다중선형성과 정규화 조건을 만족함을 증명할 수 있습니다.\item 정규화된 교대 다중선형 함수의 유일성을 설명할 수 있습니다.}{\item 순열(permutation)과 부호(sign)\item 표준기저와 다중선형성}

행렬식의 존재와 유일성을 함께 다루기 위해 먼저 부호 있는 순열 합을 정의합니다.

\begin{definition}
$S_n$을 $\{1,\dots,n\}$의 모든 순열의 집합이라 하자. $\sigma\in S_n$의 부호를 $\operatorname{sgn}(\sigma)\in\{1,-1\}$로 쓴다.
\end{definition}

\begin{definition}
$A=(a_{ij})\in M_n(K)$에 대하여
$$
\det(A)=\sum_{\sigma\in S_n}\operatorname{sgn}(\sigma)\prod_{i=1}^n a_{i,\sigma(i)}
$$
로 정의한다. 이를 Leibniz 공식에 의한 행렬식이라 한다.
\end{definition}

\begin{theorem}
위의 정의로 주어진 $\det:M_n(K)\to K$는 다음을 만족한다.
\begin{enumerate}[label=(\roman*)]
\item 각 행에 대하여 선형이다.
\item 두 행이 같으면 행렬식은 $0$이다.
\item $\det(I_n)=1$이다.
\end{enumerate}
\end{theorem}

\paragraph{증명 전략.}
Leibniz 합에서 각 항이 한 행의 성분을 정확히 한 번만 포함한다는 점을 이용해 선형성을 보이고, 두 행이 같은 경우에는 순열을 쌍으로 묶어 소거되게 하겠습니다. 마지막으로 단위행렬에서는 항 하나만 살아남음을 확인하겠습니다.

\begin{proof}
$A$의 $k$번째 행을 $u=(u_1,\dots,u_n)$로 바꾼 행렬을 $A[k\leftarrow u]$로 쓰자. 그러면
$$
\det(A[k\leftarrow u+v])
=\sum_{\sigma}\operatorname{sgn}(\sigma)\left(\prod_{i\ne k}a_{i,\sigma(i)}\right)(u_{\sigma(k)}+v_{\sigma(k)}).
$$
분배법칙으로
$$
\det(A[k\leftarrow u+v])=\det(A[k\leftarrow u])+\det(A[k\leftarrow v]).
$$
마찬가지로 임의의 $c\in K$에 대해
$$
\det(A[k\leftarrow cu])=c\det(A[k\leftarrow u]).
$$
따라서 각 행에 대한 선형성이 성립한다.

이제 $p\ne q$번째 행이 같다고 하자. 각 $\sigma\in S_n$에 대해 $\tau=\sigma\circ(p\,q)$를 두면 $\operatorname{sgn}(\tau)=-\operatorname{sgn}(\sigma)$이다. 또한 $p,q$행이 같으므로
$$
\prod_{i=1}^n a_{i,\sigma(i)}=\prod_{i=1}^n a_{i,\tau(i)}.
$$
따라서 $\sigma,\tau$에 대응하는 두 항의 합은 $0$이고, 모든 항이 이렇게 소거되므로 $\det(A)=0$이다.

마지막으로 $I_n$에서 항
$$
\prod_{i=1}^n (I_n)_{i,\sigma(i)}
$$
가 $0$이 아니려면 모든 $i$에 대해 $\sigma(i)=i$이어야 한다. 즉 $\sigma=\operatorname{id}$ 항만 살아남고 그 값은 $1$이다. 따라서 $\det(I_n)=1$.
\end{proof}

\paragraph{의미 해석.}
이 정리는 행렬식이 ``교대 다중선형 함수''라는 구조를 실제 계산식에서 직접 읽어낼 수 있음을 보여 줍니다. 이후 성질들은 이 구조에서 자동으로 나오므로, 공식 암기보다 구조 이해가 훨씬 강력합니다.

\begin{theorem}
$F:M_n(K)\to K$가 각 행에 대해 선형이고 두 행이 같으면 $0$이 되는 함수라 하자. 그러면 모든 $A\in M_n(K)$에 대하여
$$
F(A)=F(I_n)\det(A)
$$
가 성립한다. 특히 $F(I_n)=1$이면 $F=\det$이다.
\end{theorem}

\paragraph{증명 전략.}
행을 표준기저의 선형결합으로 전개하여 $F(A)$를 기저 벡터 배열에서의 값들의 합으로 바꾸겠습니다. 반복된 기저 인덱스는 교대성 때문에 사라지고, 결국 순열 항만 남아 Leibniz 공식과 일치함을 보이겠습니다.

\begin{proof}
$A=(a_{ij})$의 $i$번째 행을 $r_i=\sum_{j=1}^n a_{ij}e_j$라 두자($e_j$는 표준기저). 다중선형성으로
$$
F(A)=\sum_{j_1,\dots,j_n} a_{1j_1}\cdots a_{nj_n}\,F(e_{j_1},\dots,e_{j_n}).
$$
만약 어떤 $p\ne q$에서 $j_p=j_q$이면 교대성에 의해 $F(e_{j_1},\dots,e_{j_n})=0$이다. 따라서 살아남는 항은 $(j_1,\dots,j_n)=(\sigma(1),\dots,\sigma(n))$ 꼴뿐이다.

교대성과 다중선형성으로 인자 두 개를 맞바꾸면 부호가 바뀐다. 실제로 $u,v$에 대해
$$
0=F(\dots,u+v,\dots,u+v,\dots)
$$
를 전개하면
$$
F(\dots,u,\dots,v,\dots)+F(\dots,v,\dots,u,\dots)=0
$$
이므로
$$
F(\dots,v,\dots,u,\dots)=-F(\dots,u,\dots,v,\dots).
$$
따라서 임의의 순열 $\sigma$에 대해
$$
F(e_{\sigma(1)},\dots,e_{\sigma(n)})=\operatorname{sgn}(\sigma)F(e_1,\dots,e_n)=\operatorname{sgn}(\sigma)F(I_n).
$$
결국
$$
F(A)=F(I_n)\sum_{\sigma\in S_n}\operatorname{sgn}(\sigma)\prod_{i=1}^n a_{i,\sigma(i)}
=F(I_n)\det(A).
$$
증명 끝.
\end{proof}

\paragraph{의미 해석.}
행렬식의 유일성은 ``좋아 보이는 함수가 많아 보이지만 실제로는 하나뿐''이라는 사실을 말합니다. 이 때문에 이후 정리에서 어떤 함수가 교대 다중선형임만 확인하면 즉시 행렬식의 상수배임을 결론낼 수 있습니다.

\begin{example}
$A=\begin{bmatrix}1&2&0\\0&3&4\\5&0&6\end{bmatrix}$에 대해 Leibniz 공식을 쓰면, 영이 아닌 항은 $\sigma=\operatorname{id}$와 $\sigma=(1\,2\,3)$ 두 개뿐입니다. 따라서
$$
\det(A)=1\cdot3\cdot6+2\cdot4\cdot5=18+40=58.
$$
\end{example}

\subsection*{주의}
\begin{warning}
행렬식은 정사각행렬에서만 정의합니다. $m\times n$ 행렬($m\ne n$)에 대해 $\det$를 쓰면 문맥 오류입니다.
\end{warning}
\begin{warning}
``선형''은 행렬 전체에 대한 선형이 아니라 ``각 행을 고정하고 한 행만 바꿀 때''의 선형입니다. $\det(A+B)=\det(A)+\det(B)$는 일반적으로 성립하지 않습니다.
\end{warning}

\subsection*{자가진단퀴즈}
\begin{enumerate}[label=\arabic*.]
\item $2\times2$ 행렬에 대해 Leibniz 공식을 직접 전개하여 $\det\begin{bmatrix}a&b\\c&d\end{bmatrix}=ad-bc$를 유도하십시오.
\item 정리 1의 증명에서 ``두 행이 같으면 0''을 순열 쌍 소거 방식으로 다시 작성해 보십시오.
\item $F(I_n)=3$인 교대 다중선형 함수 $F$가 있을 때 $F(A)$를 $\det(A)$로 표현하십시오.
\end{enumerate}

\sectiontemplate{행연산, 삼각행렬, 곱셈성}{행렬식의 실전 계산은 행연산과 삼각화에 크게 의존합니다. 이 절에서는 계산 규칙의 근거를 증명하고, 핵심 성질인 곱셈성까지 연결하겠습니다.}{\item 기본 행연산이 행렬식에 미치는 영향을 증명할 수 있습니다.\item 삼각행렬의 행렬식을 즉시 계산할 수 있습니다.\item $\det(AB)=\det(A)\det(B)$와 $\det(A^T)=\det(A)$를 엄밀하게 사용할 수 있습니다.}{\item 기본 행연산\item 행렬 곱과 전치}

실제 계산 과정에서 가장 많이 쓰는 것은 행연산 규칙입니다.

\begin{definition}
다음 세 가지를 기본 행연산이라 한다.
\begin{enumerate}[label=(\roman*)]
\item 두 행 교환
\item 한 행에 스칼라배
\item 한 행에 다른 행의 스칼라배를 더함
\end{enumerate}
\end{definition}

\begin{theorem}
행렬 $A$에 기본 행연산을 적용해 $B$를 얻었다고 하자.
\begin{enumerate}[label=(\roman*)]
\item 한 행에 $c$배를 하면 $\det(B)=c\det(A)$.
\item 한 행에 다른 행의 $c$배를 더하면 $\det(B)=\det(A)$.
\item 두 행을 교환하면 $\det(B)=-\det(A)$.
\end{enumerate}
\end{theorem}

\paragraph{증명 전략.}
정리 8.1.1에서 얻은 다중선형성과 교대성을 직접 대입하겠습니다. (ii)는 선형성으로 분해한 뒤 반복행 소거를 쓰고, (iii)는 $r_i+r_j$를 두 위치에 동시에 넣어 전개하는 표준 아이디어를 사용합니다.

\begin{proof}
(i)와 (ii)는 다중선형성에서 즉시 나온다. (ii)를 쓰면, $i\ne j$에 대해
$$
\det(\dots,r_i+cr_j,\dots,r_j,\dots)
=\det(\dots,r_i,\dots,r_j,\dots)+c\det(\dots,r_j,\dots,r_j,\dots)
=\det(\dots,r_i,\dots,r_j,\dots).
$$

(iii)를 보이자. $i\ne j$에 대해
$$
0=\det(\dots,r_i+r_j,\dots,r_i+r_j,\dots)
$$
를 전개하면
$$
0=\det(\dots,r_i,\dots,r_j,\dots)+\det(\dots,r_j,\dots,r_i,\dots)
$$
가 된다(두 개의 ``같은 행'' 항은 교대성으로 $0$). 따라서
$$
\det(\dots,r_j,\dots,r_i,\dots)=-\det(\dots,r_i,\dots,r_j,\dots).
$$
즉 두 행 교환 시 부호가 반전된다.
\end{proof}

\paragraph{의미 해석.}
가우스 소거법에서 왜 어떤 연산은 값이 유지되고, 어떤 연산은 스케일 또는 부호를 바꾸는지가 이 정리로 완전히 설명됩니다. 계산 규칙은 기억 대상이 아니라 구조의 직접 결과입니다.

\begin{theorem}
$A\in M_n(K)$가 상삼각행렬이면
$$
\det(A)=\prod_{i=1}^n a_{ii}.
$$
\end{theorem}

\paragraph{증명 전략.}
Leibniz 합에서 영이 아닌 항이 언제 가능한지를 조건 $a_{ij}=0$ ($i>j$)로 분석하겠습니다. 비항등 순열은 반드시 어떤 $i$에서 $\sigma(i)<i$를 만들므로 해당 항이 소거됩니다.

\begin{proof}
상삼각행렬에서는 $i>j$이면 $a_{ij}=0$이다. Leibniz 항
$$
\prod_{i=1}^n a_{i,\sigma(i)}
$$
가 영이 아니려면 모든 $i$에 대해 $\sigma(i)\ge i$여야 한다. 그런데 순열의 값 합은
$$
\sum_{i=1}^n \sigma(i)=\sum_{i=1}^n i
$$
이므로 $\sigma(i)\ge i$가 모두 성립하면 실제로 $\sigma(i)=i$가 모든 $i$에서 성립한다. 즉 $\sigma=\operatorname{id}$만 가능하다. 따라서
$$
\det(A)=\prod_{i=1}^n a_{ii}.
$$
\end{proof}

\paragraph{의미 해석.}
삼각화가 계산에서 강력한 이유가 명확해집니다. 소거 과정으로 삼각형태를 만들면 행렬식은 대각 곱 하나로 끝납니다.

\begin{theorem}
임의의 $A,B\in M_n(K)$에 대해
$$
\det(AB)=\det(A)\det(B)
$$
가 성립한다.
\end{theorem}

\paragraph{증명 전략.}
$B$를 고정하고 $F_B(A)=\det(AB)$를 정의한 뒤, $F_B$가 교대 다중선형임을 확인하겠습니다. 그러면 정리 8.1.2의 유일성으로 $F_B(A)=F_B(I_n)\det(A)$를 즉시 얻습니다.

\begin{proof}
$B$를 고정하고
$$
F_B(A)=\det(AB)
$$
라 두자. $A$의 $i$번째 행을 $r_i$라 하면 $(AB)$의 $i$번째 행은 $r_iB$이므로 $r_i$에 대해 선형이다. 따라서 $F_B$는 각 행에 대해 선형이다. 또한 $A$의 두 행이 같으면 $AB$의 대응 두 행도 같으므로 $F_B(A)=0$. 즉 $F_B$는 교대 다중선형 함수다.

정리 8.1.2에 의해
$$
F_B(A)=F_B(I_n)\det(A).
$$
그런데
$$
F_B(I_n)=\det(I_nB)=\det(B).
$$
따라서
$$
\det(AB)=\det(B)\det(A)=\det(A)\det(B).
$$
\end{proof}

\paragraph{의미 해석.}
곱셈성은 행렬식이 ``선형변환의 스케일링 효과''를 측정한다는 직관과 정확히 일치합니다. 변환을 연속 적용하면 스케일이 곱해지는 구조가 그대로 나타납니다.

\begin{theorem}
임의의 $A\in M_n(K)$에 대해
$$
\det(A^T)=\det(A)
$$
가 성립한다.
\end{theorem}

\paragraph{증명 전략.}
Leibniz 공식에서 전치로 인해 행/열 인덱스가 바뀐 항을 $\sigma^{-1}$ 치환으로 정리하겠습니다. 순열의 부호가 역순열에서 보존된다는 사실을 사용합니다.

\begin{proof}
$A=(a_{ij})$라 하자. 그러면
$$
\det(A^T)=\sum_{\sigma\in S_n}\operatorname{sgn}(\sigma)\prod_{i=1}^n a_{\sigma(i),i}.
$$
여기서 $\tau=\sigma^{-1}$로 두면 $\operatorname{sgn}(\tau)=\operatorname{sgn}(\sigma)$이고
$$
\prod_{i=1}^n a_{\sigma(i),i}=\prod_{j=1}^n a_{j,\tau(j)}.
$$
따라서
$$
\det(A^T)=\sum_{\tau\in S_n}\operatorname{sgn}(\tau)\prod_{j=1}^n a_{j,\tau(j)}=\det(A).
$$
\end{proof}

\paragraph{의미 해석.}
행과 열의 역할이 본질적으로 대칭이라는 점이 확인됩니다. 이 결과 덕분에 ``열 기준'' 성질도 ``행 기준'' 증명에서 바로 얻을 수 있습니다.

\begin{example}
$$
A=\begin{bmatrix}
1&2&3\\
0&1&4\\
2&1&0
\end{bmatrix}
$$
에서
$$
R_3\leftarrow R_3-2R_1,\qquad
R_3\leftarrow R_3+3R_2
$$
를 적용하면
$$
\begin{bmatrix}
1&2&3\\
0&1&4\\
0&0&6
\end{bmatrix}
$$
가 된다. 두 연산은 (ii)형이므로 행렬식이 유지되고, 삼각행렬 정리에 의해
$$
\det(A)=1\cdot1\cdot6=6.
$$
\end{example}

\subsection*{주의}
\begin{warning}
소거 과정에서 한 행을 $c$배하면 행렬식도 $c$배됩니다. ``삼각화 후 대각 곱''만 계산하면 중간 스케일 계수를 놓치기 쉽습니다.
\end{warning}
\begin{warning}
$\det(AB)=\det(A)\det(B)$는 성립하지만 일반적으로 $AB=BA$는 성립하지 않습니다. 행렬식의 교환법칙과 행렬곱의 교환법칙을 혼동하면 안 됩니다.
\end{warning}

\subsection*{자가진단퀴즈}
\begin{enumerate}[label=\arabic*.]
\item 하삼각행렬의 행렬식 공식이 왜 동일한지 $\det(A^T)=\det(A)$를 이용해 설명하십시오.
\item $R_i\leftarrow cR_i$와 $R_i\leftrightarrow R_j$를 연속 적용한 뒤 행렬식 보정 계수를 어떻게 계산하는지 규칙을 정리하십시오.
\item $\det(AB)=\det(BA)$를 직접 증명해 보십시오.
\end{enumerate}

\sectiontemplate{여인수 전개와 수반행렬}{Leibniz 공식은 원리적으로 완전하지만 계산에 비효율적일 때가 많습니다. 여인수 전개(cofactor expansion)와 수반행렬(adjugate matrix)은 계산을 체계화하고 역행렬 공식까지 제공합니다.}{\item 소행렬식(minor), 여인수(cofactor), 수반행렬(adjugate)을 정확히 정의할 수 있습니다.\item Laplace 전개 공식을 엄밀히 증명할 수 있습니다.\item $A\operatorname{adj}(A)=\det(A)I$와 역행렬 공식을 유도할 수 있습니다.}{\item 전치와 열 전개 대칭성\item 가역행렬의 정의}

\begin{definition}
$A=(a_{ij})\in M_n(K)$와 $1\le i,j\le n$에 대해
\begin{enumerate}[label=(\roman*)]
\item $M_{ij}$: $i$행과 $j$열을 지운 $(n-1)\times(n-1)$ 행렬의 행렬식(소행렬식)
\item $C_{ij}=(-1)^{i+j}M_{ij}$: $a_{ij}$의 여인수
\item $\operatorname{adj}(A)=(C_{ji})$: 여인수행렬의 전치, 즉 수반행렬
\end{enumerate}
로 정의한다.
\end{definition}

\begin{theorem}[Laplace 전개]
임의의 $A\in M_n(K)$와 고정된 $i,j$에 대해
$$
\det(A)=\sum_{k=1}^n a_{ik}C_{ik}
\quad\text{(i행 전개)},
\qquad
\det(A)=\sum_{k=1}^n a_{kj}C_{kj}
\quad\text{(j열 전개)}.
$$
\end{theorem}

\paragraph{증명 전략.}
행 전개는 Leibniz 합을 $\sigma(i)=k$인 부분합으로 분할해 얻겠습니다. 부호 인자는 해당 위치의 행/열을 끝으로 이동시키는 순열을 도입해 계산하고, 열 전개는 전치 불변성으로 귀결하겠습니다.

\begin{proof}
먼저 행 전개를 보인다. $i$를 고정하고
$$
S_n(i,k)=\{\sigma\in S_n:\sigma(i)=k\}
$$
로 두면
$$
\det(A)=\sum_{k=1}^n\sum_{\sigma\in S_n(i,k)}
\operatorname{sgn}(\sigma)\prod_{r=1}^n a_{r,\sigma(r)}.
$$
고정된 $k$에 대해 $\sigma\in S_n(i,k)$에서 $i$행과 $k$열을 삭제하면 $S_{n-1}$의 순열 $\bar\sigma$가 얻어지고, 이는 일대일 대응이다.

이때 부호 관계를 계산한다. $p=(i\,i+1\,\cdots\,n)$, $q=(k\,k+1\,\cdots\,n)$라 두면
$$
\operatorname{sgn}(p)=(-1)^{n-i},\qquad
\operatorname{sgn}(q)=(-1)^{n-k}.
$$
또한 $q\sigma p^{-1}$는 $n$을 고정하고 나머지에서 $\bar\sigma$를 유도하므로
$$
\operatorname{sgn}(\bar\sigma)=\operatorname{sgn}(q\sigma p^{-1})
=\operatorname{sgn}(q)\operatorname{sgn}(\sigma)\operatorname{sgn}(p^{-1})
=(-1)^{i+k}\operatorname{sgn}(\sigma).
$$
따라서
$$
\operatorname{sgn}(\sigma)=(-1)^{i+k}\operatorname{sgn}(\bar\sigma).
$$
또한 $\sigma(i)=k$이므로
$$
\prod_{r=1}^n a_{r,\sigma(r)}
=a_{ik}\prod_{r\ne i}a_{r,\sigma(r)}.
$$
결국
$$
\sum_{\sigma\in S_n(i,k)}\operatorname{sgn}(\sigma)\prod_{r=1}^n a_{r,\sigma(r)}
=a_{ik}(-1)^{i+k}\!\!\sum_{\bar\sigma\in S_{n-1}}\operatorname{sgn}(\bar\sigma)\prod_{r\ne i}a_{r,\sigma(r)}
=a_{ik}C_{ik}.
$$
$k$에 대해 합하면
$$
\det(A)=\sum_{k=1}^n a_{ik}C_{ik}.
$$

이제 열 전개는 $A^T$에 대해 같은 식을 적용하면
$$
\det(A^T)=\sum_{k=1}^n (A^T)_{jk}\,C^{A^T}_{jk}
=\sum_{k=1}^n a_{kj}C_{kj}.
$$
정리 8.2.4의 $\det(A^T)=\det(A)$를 사용하면 열 전개가 성립한다.
\end{proof}

\paragraph{의미 해석.}
Laplace 전개는 ``한 행(또는 한 열) 기준으로 차원을 1 낮추는 재귀 계산''을 가능하게 합니다. 특히 영이 많은 행/열을 선택하면 계산량이 크게 줄어듭니다.

\begin{theorem}
임의의 $A\in M_n(K)$에 대해
$$
A\operatorname{adj}(A)=\det(A)I_n,
\qquad
\operatorname{adj}(A)A=\det(A)I_n
$$
가 성립한다.
\end{theorem}

\paragraph{증명 전략.}
왼쪽 곱의 $(i,j)$성분과 오른쪽 곱의 $(i,j)$성분을 각각 계산하겠습니다. 대각성분은 Laplace 전개로 $\det(A)$가 되고, 비대각성분은 ``한 행(또는 열) 복사'' 행렬의 행렬식이 $0$임을 이용해 소거됩니다.

\begin{proof}
먼저 $(A\operatorname{adj}(A))_{ij}$를 계산하면
$$
(A\operatorname{adj}(A))_{ij}
=\sum_{k=1}^n a_{ik}C_{jk}.
$$
$i=j$이면 Laplace의 $j$행 전개로
$$
\sum_{k=1}^n a_{jk}C_{jk}=\det(A).
$$
$i\ne j$이면 $B$를 $A$의 $j$행을 $i$행으로 바꾼 행렬이라 하자. $B$는 두 행이 같으므로 $\det(B)=0$. 이를 $j$행으로 전개하면
$$
0=\sum_{k=1}^n b_{jk}C_{jk}(B)
=\sum_{k=1}^n a_{ik}C_{jk}(A),
$$
왜냐하면 $j$행을 삭제하여 만드는 소행렬식은 $A$와 $B$에서 동일하기 때문이다. 따라서
$$
(A\operatorname{adj}(A))_{ij}=0\quad(i\ne j).
$$
즉
$$
A\operatorname{adj}(A)=\det(A)I_n.
$$

같은 방식으로
$$
(\operatorname{adj}(A)A)_{ij}
=\sum_{k=1}^n C_{ki}a_{kj}
$$
를 계산한다. $i=j$이면 $i$열 전개로 $\det(A)$를 얻는다. $i\ne j$이면 $D$를 $A$의 $i$열을 $j$열로 바꾼 행렬이라 두면 $\det(D)=0$이고, $i$열 전개로
$$
0=\sum_{k=1}^n d_{ki}C_{ki}(D)
=\sum_{k=1}^n a_{kj}C_{ki}(A)
=\sum_{k=1}^n C_{ki}a_{kj}.
$$
따라서 비대각성분이 0이므로
$$
\operatorname{adj}(A)A=\det(A)I_n.
$$
증명 끝.
\end{proof}

\paragraph{의미 해석.}
이 정리는 수반행렬이 ``역행렬의 분자'' 역할을 한다는 사실을 정확히 보여 줍니다. 한 번의 공식을 기억하는 것이 아니라, 왜 역행렬 공식이 가능한지 구조적으로 이해하게 됩니다.

\begin{theorem}
$A\in M_n(K)$에 대해 다음이 동치이다.
\begin{enumerate}[label=(\roman*)]
\item $A$는 가역이다.
\item $\det(A)\ne0$.
\end{enumerate}
또한 이 조건이 성립하면
$$
A^{-1}=\frac{1}{\det(A)}\operatorname{adj}(A).
$$
\end{theorem}

\paragraph{증명 전략.}
(i)$\Rightarrow$(ii)는 곱셈성으로 즉시 얻고, (ii)$\Rightarrow$(i)는 직전 정리의 항등식을 $\det(A)$로 나누어 역행렬 후보를 직접 구성하겠습니다.

\begin{proof}
(i)$\Rightarrow$(ii): $A^{-1}$가 존재하면
$$
1=\det(I_n)=\det(AA^{-1})=\det(A)\det(A^{-1}),
$$
따라서 $\det(A)\ne0$.

(ii)$\Rightarrow$(i): 정리 8.3.2에서
$$
A\operatorname{adj}(A)=\det(A)I_n,\qquad
\operatorname{adj}(A)A=\det(A)I_n.
$$
$\det(A)\ne0$이므로
$$
B=\frac{1}{\det(A)}\operatorname{adj}(A)
$$
를 두면
$$
AB=I_n,\qquad BA=I_n.
$$
따라서 $B=A^{-1}$이고 공식이 성립한다.
\end{proof}

\paragraph{의미 해석.}
행렬식은 단순 수치가 아니라 가역성의 완전한 판정자입니다. 이후 선형시스템 해 존재성, 고유값 계산, 최소다항식 분석에서 이 동치가 반복적으로 핵심 연결고리로 등장합니다.

\begin{example}
$$
A=\begin{bmatrix}2&1\\5&3\end{bmatrix}
$$
이면
$$
\det(A)=2\cdot3-1\cdot5=1.
$$
여인수는
$$
C_{11}=3,\quad C_{12}=-5,\quad C_{21}=-1,\quad C_{22}=2
$$
이므로
$$
\operatorname{adj}(A)=\begin{bmatrix}3&-1\\-5&2\end{bmatrix},
\qquad
A^{-1}=\operatorname{adj}(A).
$$
\end{example}

\subsection*{주의}
\begin{warning}
여인수 부호는 체스판 부호
$$
\begin{bmatrix}
+&-&+&\cdots\\
-&+&-&\cdots\\
+&-&+&\cdots
\end{bmatrix}
$$
를 따릅니다. 부호 하나만 틀려도 역행렬 전체가 틀립니다.
\end{warning}
\begin{warning}
$\operatorname{adj}(A)$는 원소별 역수 행렬이 아닙니다. ``소행렬식 계산 $\to$ 부호 $\to$ 전치''의 세 단계를 반드시 지켜야 합니다.
\end{warning}

\subsection*{자가진단퀴즈}
\begin{enumerate}[label=\arabic*.]
\item $3\times3$ 행렬 하나를 정해 특정 행 기준 Laplace 전개를 직접 계산하십시오.
\item 정리 8.3.2에서 비대각성분이 0이 되는 이유를 ``행 복사'' 행렬로 다시 설명하십시오.
\item $\det(A)=0$일 때 역행렬 공식이 왜 실패하는지 분모 관점과 선형독립 관점에서 각각 설명하십시오.
\end{enumerate}

\sectiontemplate{Cramer 공식과 특수 행렬식}{행렬식은 역행렬 판정에 그치지 않고 연립일차방정식 해 공식과 보간 문제의 가역성 판정에도 직접 쓰입니다. 이 절에서는 Cramer 공식과 Vandermonde 행렬식을 완전증명으로 정리하겠습니다.}{\item Cramer 공식을 정리와 증명으로 설명할 수 있습니다.\item Vandermonde 행렬식 공식을 유도할 수 있습니다.\item 두 공식을 연립방정식/보간 문제에 적용할 수 있습니다.}{\item 열 다중선형성\item 다항식의 근과 인수정리}

\begin{definition}
$A=[a_1\,a_2\,\cdots\,a_n]\in M_n(K)$, $b\in K^n$라 하자. $A$의 $j$번째 열을 $b$로 바꾼 행렬을
$$
A_j(b)=[a_1\,\cdots\,a_{j-1}\,b\,a_{j+1}\,\cdots\,a_n]
$$
로 정의한다.
\end{definition}

\begin{theorem}[Cramer 공식]
$\det(A)\ne0$인 $n\times n$ 연립방정식
$$
Ax=b
$$
의 유일해 $x=(x_1,\dots,x_n)^T$는
$$
x_j=\frac{\det(A_j(b))}{\det(A)}\qquad(1\le j\le n)
$$
를 만족한다.
\end{theorem}

\paragraph{증명 전략.}
해 $x$를 열벡터 조합 $b=\sum_k x_k a_k$로 쓴 뒤, $A_j(b)$의 행렬식을 $j$열 다중선형성으로 전개하겠습니다. 중복 열이 생기는 항은 0이 되어 원하는 항 하나만 남습니다.

\begin{proof}
정리 8.3.3에 의해 $\det(A)\ne0$이면 $A$는 가역이므로 해는 유일하다. 이제 $Ax=b$를 열로 쓰면
$$
b=x_1a_1+\cdots+x_na_n.
$$
또한 정리 8.2.4와 정리 8.1.1로부터 행렬식은 열에 대해서도 다중선형이다. 따라서
$$
\det(A_j(b))
=\det(a_1,\dots,a_{j-1},\sum_{k=1}^n x_ka_k,a_{j+1},\dots,a_n)
=\sum_{k=1}^n x_k\det(a_1,\dots,a_{j-1},a_k,a_{j+1},\dots,a_n).
$$
$k\ne j$이면 $k$열과 $j$열이 모두 $a_k$가 되어 두 열이 같으므로 항이 0이다. $k=j$만 남아
$$
\det(A_j(b))=x_j\det(A).
$$
$\det(A)\ne0$이므로
$$
x_j=\frac{\det(A_j(b))}{\det(A)}.
$$
\end{proof}

\paragraph{의미 해석.}
Cramer 공식은 ``해의 각 좌표를 행렬식 비율로 읽는 방법''을 줍니다. 계산량 때문에 대규모 수치해석에는 비효율적이지만, 해의 의존성 분석이나 이론 전개에서는 매우 투명한 도구입니다.

\begin{definition}
스칼라 $x_1,\dots,x_n\in K$에 대해
$$
V(x_1,\dots,x_n)=
\begin{bmatrix}
1&x_1&x_1^2&\cdots&x_1^{n-1}\\
1&x_2&x_2^2&\cdots&x_2^{n-1}\\
\vdots&\vdots&\vdots&&\vdots\\
1&x_n&x_n^2&\cdots&x_n^{n-1}
\end{bmatrix}
$$
를 Vandermonde 행렬이라 한다.
\end{definition}

\begin{theorem}[Vandermonde 행렬식]
모든 $n\ge1$에 대해
$$
\det(V(x_1,\dots,x_n))=\prod_{1\le i<j\le n}(x_j-x_i).
$$
\end{theorem}

\paragraph{증명 전략.}
$n$에 대한 귀납법을 사용합니다. 마지막 변수 $x_n$을 미지수 $t$로 보고 행렬식을 다항식 $f(t)$로 해석하면, $t=x_1,\dots,x_{n-1}$에서 영이 되어 인수 $\prod_{k=1}^{n-1}(t-x_k)$를 갖습니다. 최고차항 계수를 비교해 상수를 결정하겠습니다.

\begin{proof}
$n=1$에서는 양변이 모두 $1$이므로 자명하다.

이제 $n-1$에서 성립한다고 가정하고 $n$을 보이자. $x_1,\dots,x_{n-1}$를 고정하고
$$
f(t)=\det(V(x_1,\dots,x_{n-1},t))
$$
라 두면 $\deg f\le n-1$이다.

$k=1,\dots,n-1$에 대해 $t=x_k$이면 $k$번째 행과 마지막 행이 같아져 $f(x_k)=0$. 따라서
$$
\prod_{k=1}^{n-1}(t-x_k)\mid f(t).
$$
두 다항식의 차수가 모두 $n-1$이므로
$$
f(t)=c\prod_{k=1}^{n-1}(t-x_k)
$$
인 상수 $c$가 존재한다.

이제 $c$를 구한다. 마지막 행으로 Laplace 전개하면 $t^{n-1}$의 계수는 마지막 열 항에서만 나오고 그 계수는
$$
\det(V(x_1,\dots,x_{n-1}))
$$
이다(부호는 $(-1)^{n+n}=1$). 즉
$$
c=\det(V(x_1,\dots,x_{n-1})).
$$
귀납가정으로
$$
c=\prod_{1\le i<j\le n-1}(x_j-x_i).
$$
따라서
$$
f(t)=\left(\prod_{1\le i<j\le n-1}(x_j-x_i)\right)\left(\prod_{k=1}^{n-1}(t-x_k)\right).
$$
여기서 $t=x_n$을 대입하면
$$
\det(V(x_1,\dots,x_n))
=\prod_{1\le i<j\le n}(x_j-x_i).
$$
증명 끝.
\end{proof}

\paragraph{의미 해석.}
이 공식은 ``서로 다른 점에서의 다항식 보간이 유일하다''는 사실의 핵심 선형대수 근거입니다. 곱이 0이 아니면 Vandermonde 행렬이 가역이므로 보간계수는 하나로 정해집니다.

\begin{example}
점 $x_1=0,x_2=1,x_3=2$에 대한 Vandermonde 행렬
$$
V=
\begin{bmatrix}
1&0&0\\
1&1&1\\
1&2&4
\end{bmatrix}
$$
의 행렬식은
$$
(1-0)(2-0)(2-1)=2
$$
이므로 $0,1,2$에서의 2차 보간 문제는 항상 유일해를 갖는다.
\end{example}

\subsection*{주의}
\begin{warning}
Cramer 공식은 정사각행렬이면서 $\det(A)\ne0$일 때만 사용합니다. 직사각형 시스템이나 특이행렬에 그대로 적용하면 잘못된 결론이 나옵니다.
\end{warning}
\begin{warning}
Vandermonde 행렬은 $x_i$가 서로 매우 가까우면 수치적으로 불안정할 수 있습니다. 이론적 가역성과 수치적 안정성은 다른 개념입니다.
\end{warning}

\subsection*{자가진단퀴즈}
\begin{enumerate}[label=\arabic*.]
\item $2\times2$ 시스템 하나를 정해 Cramer 공식으로 직접 해를 구하고, 소거법 결과와 비교하십시오.
\item $\det(A_j(b))=x_j\det(A)$에서 $k\ne j$ 항이 사라지는 이유를 ``중복 열'' 관점으로 설명하십시오.
\item Vandermonde 공식에서 $x_p=x_q$가 되면 왜 행렬식이 즉시 $0$이 되는지 행의 관점으로 설명하십시오.
\end{enumerate}

\chapterapplicationtemplate{3차원 좌표 보정에서 측정 좌표 $y$와 실제 좌표 $x$가 $y=Ax$로 연결된다고 하겠습니다. 이때 $|\det(A)|$는 부피 스케일 계수이며, $\det(A)<0$이면 좌표계 방향이 반전되었음을 뜻합니다. 또한 센서 교정 단계에서 기준 점쌍으로 선형시스템을 세우면, $\det(A)\ne0$ 여부가 보정 가능성의 1차 판정 지표가 됩니다.}

\chaptersummarytemplate

\section*{종합 연습문제}

\subsection*{연습문제 A (기초 확인)}
\begin{enumerate}[label=\textbf{A\arabic*.}]
\item $\det\begin{bmatrix}3&-2\\5&4\end{bmatrix}$를 계산하시오.
\item Leibniz 공식을 사용하여 $\det\begin{bmatrix}1&2&0\\0&1&3\\4&0&2\end{bmatrix}$를 계산하시오.
\item 상삼각행렬
$$
\begin{bmatrix}
2&1&-3\\
0&-1&5\\
0&0&4
\end{bmatrix}
$$
의 행렬식을 구하시오.
\item $A$에서 $R_2\leftarrow R_2+3R_1$을 적용해 얻은 행렬을 $B$라 할 때, $\det(B)$를 $\det(A)$로 표현하시오.
\item 행렬
$$
\begin{bmatrix}
1&0&2\\
3&-1&4\\
2&5&0
\end{bmatrix}
$$
에서 $C_{23}$을 구하시오.
\item $A=\begin{bmatrix}1&2\\3&7\end{bmatrix}$에 대해 $\operatorname{adj}(A)$와 $A^{-1}$를 구하시오.
\item 다음 연립방정식을 Cramer 공식으로 풀어라:
$$
\begin{cases}
2x+y=5,\\
x-3y=-4.
\end{cases}
$$
\item $x_1=1,x_2=3,x_3=4$일 때 Vandermonde 행렬식을 계산하시오.
\end{enumerate}

\subsection*{연습문제 B (개념 연결)}
\begin{enumerate}[label=\textbf{B\arabic*.}]
\item 임의의 $\lambda\in K$에 대해 $\det(\lambda A)=\lambda^n\det(A)$를 증명하시오.
\item $A$의 두 행이 서로 비례하면 $\det(A)=0$임을 증명하시오.
\item $A$가 가역일 때 $\det(A^{-1})=\det(A)^{-1}$임을 증명하시오.
\item 블록 상삼각행렬
$$
\begin{bmatrix}
B&*\\
0&C
\end{bmatrix}
$$
의 행렬식이 $\det(B)\det(C)$임을 증명하시오.
\item $p_A(t)=\det(tI_n-A)$가 $n$차 단위다항식임을 증명하시오.
\end{enumerate}

\subsection*{연습문제 C (증명/종합)}
\begin{enumerate}[label=\textbf{C\arabic*.}]
\item 교대 다중선형 함수 $F$가 $F(I_n)=0$이면 $F\equiv0$임을 정리 8.1.2를 사용하지 않고 직접 증명하시오.
\item $A,B$가 가역행렬일 때
$$
\operatorname{adj}(AB)=\operatorname{adj}(B)\operatorname{adj}(A)
$$
를 증명하시오.
\item 행렬 $A$의 랭크가 $r$이면 크기 $r+1$ 이상의 모든 소행렬식이 $0$임을 증명하고, 역으로 어떤 크기 $r$ 소행렬식이 0이 아니면 $\operatorname{rank}(A)\ge r$임을 증명하시오.
\end{enumerate}